% Wersja źródłowa: audyt II, jednostka 14
\setcounter{section}{13}
\section{Pochodne Gâteaux i Frécheta w przestrzeniach Banacha i Hilberta}
\label{sec:gateaux-frechet14}

\subsection{Przygotowanie: kierunki, przyrosty i liniowe przybliżenia}

W analizie funkcji jednej zmiennej pochodna opisuje najlepsze liniowe przybliżenie funkcji. Dla odwzorowania między przestrzeniami unormowanymi możliwe są dwie naturalne interpretacje: przybliżenie \emph{osobno w każdym ustalonym kierunku} oraz przybliżenie \emph{jednostajne względem wszystkich małych przyrostów}. Pierwsza prowadzi do pochodnej Gâteaux, druga do pochodnej Frécheta. Różnica jest istotna już w $\R^2$.

Niech $E,F$ będą rzeczywistymi przestrzeniami Banacha, $U\subseteq E$ zbiorem otwartym i $f:U\to F$. Przez $\mathcal L(E,F)$ oznaczamy przestrzeń ograniczonych operatorów liniowych, wyposażoną w normę operatorową. Otwartość $U$ zapewnia, że dla $x\in U$ i $h\in E$ punkty $x+th$ należą do $U$ dla wszystkich dostatecznie małych $|t|$.

\begin{definicja}[Pochodna kierunkowa]
\label{def:directional14}
Dla $x\in U$ i $h\in E$ mówimy, że $f$ ma \emph{dwustronną pochodną kierunkową} w punkcie $x$ w kierunku $h$, jeśli w $F$ istnieje granica
$\partial_h f(x)=\lim_{t\to0,\,t\in\R}(f(x+th)-f(x))/t$.
Nie wymagamy, aby zależność $h\mapsto\partial_h f(x)$ była liniowa ani ciągła. Istnienie wszystkich pochodnych kierunkowych samo w sobie nie oznacza różniczkowalności Gâteaux.
\end{definicja}

\begin{definicja}[Różniczkowalność Gâteaux]
\label{def:gateaux14}
Odwzorowanie $f$ jest \emph{różniczkowalne w sensie Gâteaux} w $x\in U$, jeśli istnieje $A\in\mathcal L(E,F)$ takie, że dla każdego $h\in E$ zachodzi $\partial_h f(x)=Ah$. Operator $A$, jednoznacznie wyznaczony przez te granice, oznaczamy $D_Gf(x)$. Przyjęliśmy definicję wymagającą ciągłości pochodnej jako operatora liniowego; w niektórych źródłach spotyka się słabszą konwencję, dopuszczającą nieograniczone odwzorowanie liniowe.
\end{definicja}

\begin{definicja}[Różniczkowalność Frécheta]
\label{def:frechet14}
Odwzorowanie $f$ jest \emph{różniczkowalne w sensie Frécheta} w $x\in U$, jeśli istnieje $A\in\mathcal L(E,F)$ oraz funkcja $r$ określona dla małych $h\in E$ taka, że
\[
f(x+h)=f(x)+Ah+r(h),\qquad
\lim_{\|h\|_E\to0,\,h\ne0}\frac{\|r(h)\|_F}{\|h\|_E}=0.
\]
Operator $A$ oznaczamy $Df(x)$ albo $D_Ff(x)$; piszemy również $r(h)=o(\|h\|_E)$. Symbol $Df(x)h$ oznacza wartość operatora $Df(x)$ na wektorze $h$, nie jego iloczyn skalarny z $h$.
\end{definicja}

\begin{uwaga}[Założenia o zupełności]
\label{rem:completeness14}
Do samych definicji wystarczą przestrzenie unormowane i otwarta dziedzina. Zupełność potrzebna będzie między innymi do twierdzenia o całkowaniu funkcji o wartościach w przestrzeni Banacha oraz w późniejszych twierdzeniach o funkcji odwrotnej. Stąd w niniejszym rozdziale pracujemy przede wszystkim w przestrzeniach Banacha.
\end{uwaga}

\subsection{Implikacje między pochodnymi}

\begin{twierdzenie}[Fréchet implikuje Gâteaux]
\label{thm:frechet-implies-gateaux14}
Niech $E,F$ będą rzeczywistymi przestrzeniami unormowanymi, $U\subseteq E$ zbiorem otwartym, a $f:U\to F$ odwzorowaniem różniczkowalnym w sensie Frécheta w punkcie $x\in U$. Wówczas $f$ jest ciągłe w $x$, jest różniczkowalne w sensie Gâteaux w $x$ oraz $D_Gf(x)=Df(x)$. Pochodna Frécheta, jeśli istnieje, jest jednoznaczna.
\end{twierdzenie}
\begin{proof}
Niech $A=Df(x)$ i $f(x+h)-f(x)=Ah+r(h)$, gdzie $\|r(h)\|=o(\|h\|)$. Ponieważ $A$ jest ograniczony, $\|Ah+r(h)\|\le\|A\|\|h\|+o(\|h\|)\to0$, co daje ciągłość. Dla ustalonego $h\ne0$ i małego rzeczywistego $t$ mamy
$\|(f(x+th)-f(x))/t-Ah\|=\|r(th)\|/|t|=\|h\|\,\|r(th)\|/\|th\|\to0$.
Dla $h=0$ pochodna kierunkowa wynosi zero. Jeżeli $A$ i $B$ spełniają definicję Frécheta, powyższy argument daje $Ah=Bh$ dla każdego $h$, a więc $A=B$.
\end{proof}

\begin{stwierdzenie}[Kryterium jednostajnego przybliżenia kierunkowego]
\label{prop:uniform-directions14}
Niech $E,F$ będą rzeczywistymi przestrzeniami unormowanymi, $U\subseteq E$ otwarty, $x\in U$ i $A\in\mathcal L(E,F)$. Odwzorowanie $f:U\to F$ ma pochodną Frécheta $Df(x)=A$ wtedy i tylko wtedy, gdy
\[
\lim_{t\to0,\,t\in\R}\ \sup_{\|h\|_E=1}
\left\|\frac{f(x+th)-f(x)}{t}-Ah\right\|_F=0.
\]
W wyrażeniu supremum rozpatrujemy wyłącznie $t\ne0$ o module na tyle małym, że $x+tB_E\subseteq U$; zakładamy też $E\ne\{0\}$.
\end{stwierdzenie}
\begin{proof}
Dla reszty $r(k)=f(x+k)-f(x)-Ak$ mamy
$\sup_{\|h\|=1}\|r(th)\|/|t|=\sup_{\|k\|=|t|}\|r(k)\|/\|k\|$.
Warunek Frécheta zapewnia, że te suprema dążą do zera. Odwrotnie, dla $k\ne0$ połóżmy $t=\|k\|$ i $h=k/\|k\|$; uzyskana jednostajna granica daje $\|r(k)\|/\|k\|\to0$.
\end{proof}

\begin{uwaga}
\label{rem:quantifiers14}
W definicji Gâteaux kwantyfikatory mają kolejność: \emph{dla każdego ustalonego $h$ granica przy $t\to0$}. W definicji Frécheta wymagamy wspólnej kontroli błędu dla wszystkich kierunków. Zamiana kolejności kwantyfikatorów nie jest uprawniona, nawet w skończonym wymiarze.
\end{uwaga}

\subsection{Kontrprzykłady: czego pochodna Gâteaux nie gwarantuje}

\begin{przyklad}[Wszystkie pochodne kierunkowe i brak ciągłości]
\label{ex:directional-discontinuous14}
Dla $(x,y)\ne(0,0)$ połóżmy
$g(x,y)=x^3y/(x^6+y^2)$, zaś $g(0,0)=0$. Dla każdego $(a,b)\in\R^2$ granica $\partial_{(a,b)}g(0,0)$ istnieje i jest równa $0$: jeśli $b\ne0$, to $g(ta,tb)/t=a^3b\,t/(a^6t^4+b^2)\to0$; jeśli $b=0$, iloraz jest zerowy. Zatem $D_Gg(0,0)=0\in\mathcal L(\R^2,\R)$. Jednak na krzywej $y=x^3$, $x\ne0$, mamy $g(x,x^3)=1/2$. Funkcja nie jest ciągła w $(0,0)$, więc nie jest tam różniczkowalna w sensie Frécheta.
\end{przyklad}

\begin{przyklad}[Ciągłość i pochodna Gâteaux nie implikują pochodnej Frécheta]
\label{ex:continuous-gateaux-not-frechet14}
Niech $f(0,0)=0$, zaś dla $(x,y)\ne(0,0)$ niech
\[
f(x,y)=\frac{x^3y}{x^4+y^2}.
\]
Z nierówności $2|x^2y|\le x^4+y^2$ wynika $|f(x,y)|\le|x|/2$; zatem $f$ jest ciągła w $(0,0)$. Dla ustalonego $(a,b)$, gdy $b\ne0$, otrzymujemy
$f(ta,tb)/t=a^3b\,t/(a^4t^2+b^2)\to0$, a dla $b=0$ iloraz jest zerowy. Stąd $D_Gf(0,0)=0$. Natomiast $f(x,x^2)=x/2$ dla $x\ne0$ oraz
$|f(x,x^2)|/\sqrt{x^2+x^4}=1/(2\sqrt{1+x^2})\to1/2$ przy $x\to0^+$. Reszta względem zerowego operatora nie jest $o(\|(x,y)\|_2)$: pochodna Frécheta nie istnieje.
\end{przyklad}

\begin{przyklad}[Sama pochodna kierunkowa nie musi być liniowa]
\label{ex:nonlinear-directional14}
Dla $f(x,y)=|x|+|y|$ w zerze istnieje tylko jednostronna pochodna kierunkowa, równa $|a|+|b|$ w kierunku $(a,b)$; dwustronna pochodna kierunkowa przy $h\ne0$ nie istnieje. Dla przykładu z \emph{dwustronnymi} pochodnymi kierunkowymi, lecz nieliniową zależnością od kierunku, rozważmy $q(0,0)=0$ oraz $q(x,y)=x^3/(x^2+y^2)$ dla $(x,y)\ne0$. Wówczas $\partial_{(a,b)}q(0,0)=a^3/(a^2+b^2)$ dla $(a,b)\ne(0,0)$, a w kierunku zerowym pochodna wynosi $0$. Odwzorowanie to jest jednorodne względem mnożenia kierunku przez rzeczywiste skalary, lecz nie jest addytywne: $\partial_{(1,0)}q+\partial_{(0,1)}q=1$, podczas gdy $\partial_{(1,1)}q=1/2$. Dlatego $q$ nie ma pochodnej Gâteaux według Definicji~\ref{def:gateaux14}.
\end{przyklad}

\begin{stwierdzenie}[Brak reguły łańcuchowej dla samej pochodnej Gâteaux]
\label{prop:chain-fails14}
Istnieją ciągłe odwzorowanie $f:\R^2\to\R$, różniczkowalne w sensie Gâteaux w zerze, oraz gładka krzywa $\gamma:\R\to\R^2$ z $\gamma(0)=0$ takie, że $\frac{\mathrm d}{\mathrm dt}(f\circ\gamma)(0)\ne D_Gf(0)\gamma'(0)$.
\end{stwierdzenie}
\begin{proof}
Niech $f$ będzie funkcją z Przykładu~\ref{ex:continuous-gateaux-not-frechet14}, a $\gamma(t)=(t,t^2)$. Wówczas $f(\gamma(t))=t/2$ dla $t\ne0$, a także dla $t=0$ z ciągłości. Zatem $(f\circ\gamma)'(0)=1/2$, podczas gdy $D_Gf(0)=0$. Niezgodność wynika z niestabilności kierunkowego przybliżenia przy zmiennym kierunku przyrostu.
\end{proof}

\begin{center}
\begin{tikzpicture}[x=2.15cm,y=1.65cm,>=Stealth]
  \draw[->,thin] (-0.17,0)--(1.7,0) node[below] {$x$};
  \draw[->,thin] (0,-0.42)--(0,1.48) node[left] {$y$};
  \draw[gray,densely dashed] (0,0)--(1.32,0.62);
  \draw[gray,densely dashed] (0,0)--(0.79,1.22);
  \draw[gray,densely dashed] (0,0)--(1.42,-0.29);
  \draw[very thick,blue!55!black,domain=0:1.15,samples=70,smooth]
       plot (\x,{\x*\x});
  \fill (0,0) circle (1.05pt);
  \node[anchor=west,blue!55!black] at (1.17,1.22) {\footnotesize $y=x^2$};
  \node[anchor=west] at (0.83,-0.37) {\scriptsize kierunki stałe};
\end{tikzpicture}

{\small Proste przez zero odpowiadają ustalonym kierunkom w definicji
Gâteaux. Krzywa $y=x^2$ prowadzi do zera po zmieniającym się
kierunku; na niej ujawnia się brak różniczkowalności Frécheta.}
\end{center}

\subsection{Pochodna Gâteaux normy w \texorpdfstring{$\ell^1$}{ell1} bez pochodnej Frécheta}

\begin{twierdzenie}[Kontrprzykład w nieskończonym wymiarze]
\label{thm:l1-gateaux14}
Niech $E=\ell^1(\N;\R)$ i $N:E\to\R$, $N(x)=\|x\|_1$. Jeśli $x=(x_n)\in E$ spełnia $x_n\ne0$ dla wszystkich $n$, to $N$ jest różniczkowalne w sensie Gâteaux w $x$ i
$D_GN(x)h=\sum_{n=1}^\infty \operatorname{sgn}(x_n)h_n$.
Funkcja $N$ nie jest jednak różniczkowalna w sensie Frécheta w żadnym takim punkcie $x$.
\end{twierdzenie}
\begin{proof}
Dla ustalonego $h\in\ell^1$ i $t\ne0$ oznaczmy $a_n(t)=(|x_n+th_n|-|x_n|)/t$. Ponieważ $x_n\ne0$, mamy $a_n(t)\to\operatorname{sgn}(x_n)h_n$ dla każdego $n$, przy czym $|a_n(t)|\le|h_n|$. Zbieżność zmajoryzowana dla szeregu (albo argument z obcięciem sumy i oszacowaniem ogona przez $\sum_{n>N}|h_n|$) daje
$\lim_{t\to0}(N(x+th)-N(x))/t=\sum_n \operatorname{sgn}(x_n)h_n$.
Jest to funkcjonał liniowy i ciągły o normie $1$.

Jeżeli istniałaby pochodna Frécheta, musiałaby być równa $D_GN(x)$. Dla $h^{(n)}=-2x_ne_n$ mamy $\|h^{(n)}\|_1=2|x_n|\to0$, ponieważ $x\in\ell^1$, a zarazem $N(x+h^{(n)})=N(x)$ i $D_GN(x)h^{(n)}=-2|x_n|$. Wobec tego
\[
\frac{|N(x+h^{(n)})-N(x)-D_GN(x)h^{(n)}|}{\|h^{(n)}\|_1}=1
\]
dla każdego $n$. Jest to sprzeczne z warunkiem Frécheta.
\end{proof}

\begin{uwaga}
Jeśli pewna współrzędna $x_n$ znika, norma $\ell^1$ nie ma w tym punkcie nawet dwustronnej pochodnej kierunkowej w kierunku $e_n$. Powyższy kontrprzykład pokazuje więc ścisłą różnicę między Gâteaux a Fréchetem również w punktach, w których wszystkie pochodne kierunkowe tworzą ciągły funkcjonał liniowy.
\end{uwaga}

\subsection{Reguły różniczkowania w przestrzeniach Banacha}

\begin{stwierdzenie}[Suma i operator liniowy]
\label{prop:sum-linear14}
Niech $E,F,G$ będą rzeczywistymi przestrzeniami Banacha, $U\subseteq E$ otwarty i $x\in U$. Jeżeli $f,g:U\to F$ są różniczkowalne w sensie Frécheta w $x$, a $\alpha,\beta\in\R$, to $D(\alpha f+\beta g)(x)=\alpha Df(x)+\beta Dg(x)$. Jeśli $T\in\mathcal L(F,G)$, to $D(T\circ f)(x)=T\circ Df(x)$. Ponadto dla $L\in\mathcal L(E,F)$ i $b\in F$ funkcja $f(y)=Ly+b$ ma $Df(y)=L$ w każdym $y\in E$.
\end{stwierdzenie}
\begin{proof}
Wszystkie równości wynikają przez podstawienie rozwinięć $f(x+h)=f(x)+Df(x)h+r_f(h)$ oraz $g(x+h)=g(x)+Dg(x)h+r_g(h)$ i użycie $\|Tr_f(h)\|\le\|T\|\|r_f(h)\|$. W przypadku funkcji afinicznej reszta wynosi zero.
\end{proof}

\begin{twierdzenie}[Reguła łańcuchowa Frécheta]
\label{thm:chain14}
Niech $E,F,G$ będą rzeczywistymi przestrzeniami Banacha, $U\subseteq E$, $V\subseteq F$ otwarte, $f:U\to F$ oraz $g:V\to G$ i $f(U)\subseteq V$. Jeśli $f$ jest różniczkowalne w sensie Frécheta w $x\in U$, a $g$ w $f(x)$, to $g\circ f$ jest różniczkowalne w sensie Frécheta w $x$, a
$D(g\circ f)(x)=Dg(f(x))\circ Df(x)$.
\end{twierdzenie}
\begin{proof}
Oznaczmy $A=Df(x)$, $B=Dg(f(x))$. Wówczas $f(x+h)=f(x)+Ah+r(h)$, $\|r(h)\|=o(\|h\|)$, i dla $k\to0$ mamy $g(f(x)+k)=g(f(x))+Bk+s(k)$, gdzie $\|s(k)\|=o(\|k\|)$. Ponieważ $f$ jest ciągłe w $x$, dla małego $h$ można przyjąć $k=Ah+r(h)$, a $\|k\|\le(\|A\|+1)\|h\|$. Otrzymujemy
$g(f(x+h))-g(f(x))-BAh=Br(h)+s(Ah+r(h))$.
Pierwszy składnik jest $o(\|h\|)$ przez ograniczoność $B$. Drugi także: jeśli $k\ne0$, to
$\|s(k)\|/\|h\|=(\|s(k)\|/\|k\|)(\|k\|/\|h\|)\to0$;
jeśli $k=0$, to $s(k)=0$. Dowód nie wymaga ciągłości pochodnych w otoczeniu rozważanych punktów.
\end{proof}

\begin{twierdzenie}[Reguła dla odwzorowania dwuliniowego]
\label{thm:bilinear14}
Niech $E_1,E_2,G$ będą rzeczywistymi przestrzeniami Banacha i $B:E_1\times E_2\to G$ ciągłym odwzorowaniem dwuliniowym, tj. $\|B(u,v)\|\le C\|u\|\|v\|$. Wówczas $B$ jest różniczkowalne w sensie Frécheta na przestrzeni produktowej $E_1\times E_2$ i
$DB(u,v)(h,k)=B(h,v)+B(u,k)$.
Jeżeli $f:U\to E_1$, $g:U\to E_2$ są różniczkowalne w $x\in U\subseteq E$ otwartym, to
$D(B(f,g))(x)h=B(Df(x)h,g(x))+B(f(x),Dg(x)h)$.
\end{twierdzenie}
\begin{proof}
Dokładna tożsamość
$B(u+h,v+k)-B(u,v)-B(h,v)-B(u,k)=B(h,k)$
daje $\|B(h,k)\|\le C\|h\|\|k\|\le C\max(\|h\|,\|k\|)^2$. Reszta jest więc $o(\max(\|h\|,\|k\|))$, a operator liniowy wskazany w tezie jest ograniczony. Ostatni wzór wynika z reguły łańcuchowej zastosowanej do odwzorowania $z\mapsto(f(z),g(z))$ i do $B$.
\end{proof}

\begin{przyklad}[Operatory i funkcje kwadratowe]
\label{ex:quadratic14}
Niech $E$ będzie rzeczywistą przestrzenią Banacha, $F$ rzeczywistą przestrzenią Banacha i $B:E\times E\to F$ ciągłym odwzorowaniem dwuliniowym. Dla $Q(x)=B(x,x)$ zachodzi
$DQ(x)h=B(h,x)+B(x,h)$.
W szczególności, jeśli $H$ jest rzeczywistą przestrzenią Hilberta i $\Phi(x)=\frac12\|x\|^2$, to
$D\Phi(x)h=\langle h,x\rangle$, a reszta ma postać $\frac12\|h\|^2$.
\end{przyklad}

\subsection{Różniczkowalność klasy \texorpdfstring{$C^1$}{C1} i twierdzenie o przyrostach}

\begin{definicja}[Klasa $C^1$]
\label{def:c1-14}
Niech $E,F$ będą rzeczywistymi przestrzeniami Banacha, $U\subseteq E$ otwarty i $f:U\to F$. Piszemy $f\in C^1(U,F)$, jeśli $f$ jest różniczkowalne w sensie Frécheta w każdym punkcie $U$ oraz odwzorowanie $Df:U\to\mathcal L(E,F)$ jest ciągłe w normie operatorowej.
\end{definicja}

\begin{lemat}[Podstawowe twierdzenie rachunku dla krzywych Banacha]
\label{lem:ftc-banach14}
Niech $F$ będzie rzeczywistą przestrzenią Banacha, $a<b$ i $\gamma:[a,b]\to F$ funkcją klasy $C^1$ (pochodna w końcach jednostronna i ciągle przedłużona). Wówczas $\gamma'$ jest całkowalna w sensie Riemanna z wartościami w $F$ oraz
$\gamma(b)-\gamma(a)=\int_a^b\gamma'(t)\,\mathrm dt$.
\end{lemat}
\begin{proof}
Ciągłość $\gamma'$ na zwartym przedziale daje jej jednostajną ciągłość, a przez kryterium Cauchy'ego sum Riemanna i zupełność $F$ zapewnia istnienie całki. Dla podziału $a=t_0<\cdots<t_m=b$ wybierzmy dowolne $\xi_j\in[t_{j-1},t_j]$. Z definicji pochodnej i jednostajnej ciągłości $\gamma'$ wynika, że przy dostatecznie małej średnicy podziału
$\|\gamma(t_j)-\gamma(t_{j-1})-(t_j-t_{j-1})\gamma'(\xi_j)\|\le\varepsilon(t_j-t_{j-1})$.
Ściślej, dla dowolnego rzeczywistego funkcjonału $\ell\in F'$ o normie co najwyżej $1$ stosujemy skalarne twierdzenie o wartości średniej do $t\mapsto\ell(\gamma(t)-t\gamma'(\xi_j))$ na $[t_{j-1},t_j]$. Otrzymujemy górne oszacowanie modułu przyrostu przez $(t_j-t_{j-1})\sup_{t\in[t_{j-1},t_j]}\|\gamma'(t)-\gamma'(\xi_j)\|$. Z twierdzenia Hahna--Banacha, po supremum po takich $\ell$, dostajemy to samo oszacowanie normy wektora. Jednostajna ciągłość $\gamma'$ zapewnia żądaną nierówność dla dostatecznie drobnych podziałów. Sumowanie oraz przejście z średnicą podziału do zera kończy dowód.
\end{proof}

\begin{twierdzenie}[Wzór całkowy i oszacowanie przyrostu]
\label{thm:mean-value14}
Niech $E,F$ będą rzeczywistymi przestrzeniami Banacha, $U\subseteq E$ zbiorem otwartym, $f\in C^1(U,F)$, a $x,y\in U$ takimi punktami, że $[x,y]=\{x+t(y-x):t\in[0,1]\}\subseteq U$. Wówczas
\[
f(y)-f(x)=\int_0^1 Df(x+t(y-x))(y-x)\,\mathrm dt,
\]
gdzie całka ma wartości w $F$. W szczególności, jeśli $\sup_{z\in[x,y]}\|Df(z)\|\le M$, to $\|f(y)-f(x)\|\le M\|y-x\|$.
\end{twierdzenie}
\begin{proof}
Definiujemy $\gamma(t)=f(x+t(y-x))$. Reguła łańcuchowa daje $\gamma'(t)=Df(x+t(y-x))(y-x)$; ciągłość $Df$ zapewnia ciągłość $\gamma'$. Lemat~\ref{lem:ftc-banach14} daje wzór. Oszacowanie normy całki przez całkę normy wynika z oszacowania sum Riemanna i przejścia do granicy.
\end{proof}

\begin{twierdzenie}[Ciągła pochodna Gâteaux daje pochodną Frécheta]
\label{thm:continuous-gateaux14}
Niech $E,F$ będą rzeczywistymi przestrzeniami Banacha, $U\subseteq E$ otwarty, $f:U\to F$. Załóżmy, że $f$ jest różniczkowalne w sensie Gâteaux w każdym punkcie $U$ i odwzorowanie $D_Gf:U\to\mathcal L(E,F)$ jest ciągłe w normie operatorowej. Wtedy $f\in C^1(U,F)$ oraz $Df(x)=D_Gf(x)$ dla wszystkich $x\in U$.
\end{twierdzenie}
\begin{proof}
Ustalmy $x\in U$ i weźmy małe $h$ takie, że $[x,x+h]\subseteq U$. Funkcja $\gamma(t)=f(x+th)$ jest różniczkowalna dla $t\in(0,1)$ i $\gamma'(t)=D_Gf(x+th)h$, co wynika z definicji kierunkowej. Jej pochodna jest ciągła dzięki założeniu, a $\gamma$ jest ciągła również w końcach przedziału, ponieważ istnienie pochodnej kierunkowej w tych punktach zapewnia ciągłość wzdłuż tej prostej. Z Lematu~\ref{lem:ftc-banach14} mamy
$f(x+h)-f(x)=\int_0^1 D_Gf(x+th)h\,\mathrm dt$.
Po odjęciu $D_Gf(x)h$ i oszacowaniu normy dostajemy
\[
\frac{\|f(x+h)-f(x)-D_Gf(x)h\|}{\|h\|}
\le\sup_{0\le t\le1}\|D_Gf(x+th)-D_Gf(x)\|\longrightarrow0.
\]
To dowodzi różniczkowalności Frécheta; ciągłość pochodnej otrzymujemy z założenia.
\end{proof}

\begin{przyklad}[Pochodna istniejąca, ale nieciągła]
\label{ex:derivative-noncontinuous14}
Na $\R$ niech $f(0)=0$ oraz $f(x)=x^2\sin(1/x^2)$ dla $x\ne0$. Wtedy $f'(0)=0$, natomiast dla $x\ne0$ mamy $f'(x)=2x\sin(1/x^2)-2x^{-1}\cos(1/x^2)$. Zatem $f$ jest różniczkowalna Frécheta w każdym punkcie, lecz nie należy do $C^1(\R,\R)$. Nie należy utożsamiać istnienia pochodnej z jej ciągłością.
\end{przyklad}

\subsection{Norma i gradient w rzeczywistych i zespolonych przestrzeniach Hilberta}

\begin{twierdzenie}[Różniczkowanie kwadratu normy i normy]
\label{thm:hilbert-norm14}
Niech $H$ będzie przestrzenią Hilberta nad $\K\in\{\R,\C\}$, traktowaną jako \emph{rzeczywista} przestrzeń Banacha, i niech $\Phi(x)=\frac12\|x\|^2$. W każdym $x\in H$ istnieje rzeczywista pochodna Frécheta $D_{\R}\Phi(x)h=\operatorname{Re}\langle h,x\rangle$. Dla $x\ne0$ norma $N(x)=\|x\|$ ma rzeczywistą pochodną Frécheta
$D_{\R}N(x)h=\operatorname{Re}\langle h,x\rangle/\|x\|$.
Jeżeli $H\ne\{0\}$, norma nie jest różniczkowalna Gâteaux w $0$ jako funkcja na rzeczywistej przestrzeni Banacha.
\end{twierdzenie}
\begin{proof}
Z rozwinięcia iloczynu skalarnego wynika dokładna tożsamość
$\Phi(x+h)-\Phi(x)-\operatorname{Re}\langle h,x\rangle=\frac12\|h\|^2$, a funkcjonał $h\mapsto\operatorname{Re}\langle h,x\rangle$ jest rzeczywisto-liniowy i ograniczony. To dowodzi pierwszej tezy. Dla $x\ne0$ stosujemy rzeczywistą regułę łańcuchową do $N(x)=\sqrt{2\Phi(x)}$, ponieważ pierwiastek jest gładki na $(0,\infty)$. W zerze iloraz kierunkowy wynosi $\|th\|/t=(|t|/t)\|h\|$; dla $h\ne0$ nie istnieje granica dwustronna.
\end{proof}

\begin{definicja}[Gradient na rzeczywistej przestrzeni Hilberta]
\label{def:gradient14}
Niech $H$ będzie rzeczywistą przestrzenią Hilberta, $U\subseteq H$ otwarty i $f:U\to\R$ różniczkowalna w sensie Frécheta w $x\in U$. \emph{Gradientem} $f$ w $x$ nazywamy jedyny wektor $\nabla f(x)\in H$ spełniający $Df(x)h=\langle h,\nabla f(x)\rangle$ dla wszystkich $h\in H$. Istnienie i jednoznaczność daje twierdzenie Riesza. Dla zespolonej przestrzeni Hilberta i funkcjonału rzeczywistego używamy struktury rzeczywistej i iloczynu skalarnego $(h,k)_{\R}=\operatorname{Re}\langle h,k\rangle$.
\end{definicja}

\begin{przyklad}[Gradient energii kwadratowej]
\label{ex:quadratic-energy14}
Niech $H$ będzie rzeczywistą lub zespoloną przestrzenią Hilberta, $A\in\mathcal L(H)$ operatorem samosprzężonym, $b\in H$, a $J(x)=\frac12\langle Ax,x\rangle-\operatorname{Re}\langle x,b\rangle$. Z samosprzężoności wynika, że $J$ ma wartości rzeczywiste. Dla $h\in H$ otrzymujemy
$D_{\R}J(x)h=\operatorname{Re}\langle h,Ax-b\rangle$, więc $\nabla J(x)=Ax-b$ w rzeczywistej strukturze Hilberta. Jeżeli $A$ jest dodatnio określony w sensie $\langle Ax,x\rangle\ge\alpha\|x\|^2$ dla pewnego $\alpha>0$, to $J$ ma dokładnie jeden minimizator: rozwiązanie $Ax=b$; jego istnienie zapewnia przypadek samosprzężony twierdzenia Laxa--Milgrama.
\end{przyklad}

\begin{przyklad}[Funkcjonał całkowy i pochodna na $L^2$]
\label{ex:integral-functional14}
Niech $(\Omega,\mathcal A,\mu)$ będzie przestrzenią z miarą skończoną i $H=L^2(\Omega,\mu;\R)$. Dla $g\in H$ określmy $J(u)=\frac12\int_\Omega |u(t)-g(t)|^2\,\mathrm d\mu(t)$. Wtedy
$DJ(u)h=\int_\Omega (u(t)-g(t))h(t)\,\mathrm d\mu(t)$,
a $\nabla J(u)=u-g$. Wynik można otrzymać przez rozwinięcie kwadratu oraz nierówność Cauchy'ego--Schwarza; reszta ma dokładnie postać $\frac12\|h\|_2^2$.
\end{przyklad}

\subsection{Rzeczywista a zespolona różniczkowalność}

\begin{definicja}[Pochodna zespolona Frécheta i Gâteaux]
\label{def:complex-differentiability14}
Niech $E,F$ będą zespolonymi przestrzeniami Banacha, $U\subseteq E$ zbiorem otwartym i $f:U\to F$. \emph{Zespoloną pochodną Frécheta} $D_{\C}f(x)$ nazywamy operator $A\in\mathcal L_{\C}(E,F)$ spełniający $f(x+h)=f(x)+Ah+o(\|h\|)$. \emph{Zespolona pochodna Gâteaux} w $x$ oznacza istnienie operatora $A\in\mathcal L_{\C}(E,F)$ takiego, że dla każdego $h\in E$
$\lim_{z\to0,\,z\in\C}(f(x+zh)-f(x))/z=Ah$.
W obu definicjach liniowość wymagana jest nad $\C$, a dla pochodnej Gâteaux parametr $z$ przebiega płaszczyznę zespoloną, nie tylko oś rzeczywistą.
\end{definicja}

\begin{stwierdzenie}[Pochodna zespolona a rzeczywista]
\label{prop:complex-real14}
Niech $E,F$ będą zespolonymi przestrzeniami Banacha, $U\subseteq E$ otwarty i $f:U\to F$. Jeśli $f$ jest zespolenie różniczkowalne w sensie Frécheta w $x\in U$, to jest też rzeczywiście różniczkowalne w sensie Frécheta w $x$, a $D_{\R}f(x)$ jest operatorem rzeczywistym odpowiadającym $D_{\C}f(x)$. Odwrotnie, jeśli $D_{\R}f(x)$ istnieje, to $D_{\C}f(x)$ istnieje wtedy i tylko wtedy, gdy $D_{\R}f(x)(ih)=iD_{\R}f(x)h$ dla wszystkich $h\in E$.
\end{stwierdzenie}
\begin{proof}
Pochodna zespolona spełnia oszacowanie reszty wymagane dla pochodnej rzeczywistej, ponieważ normy nie zmieniają się po zapomnieniu struktury zespolonej. Jeśli istnieje $D_{\R}f(x)=A$, to identyczna reszta $o(\|h\|)$ daje zespoloną pochodną dokładnie wtedy, gdy rzeczywisto-liniowy operator $A$ jest zespolenie liniowy. Dla operatora rzeczywisto-liniowego $A$ warunek $A(ih)=iA(h)$ jest równoważny liniowości zespolonej.
\end{proof}

\begin{przyklad}[Sprzężenie zespolone i kwadrat modułu]
\label{ex:complex-nonholomorphic14}
Funkcja $f:\C\to\C$, $f(z)=\overline z$, jest rzeczywiście różniczkowalna Frécheta w każdym punkcie, z $D_{\R}f(z)h=\overline h$, ale nie jest zespolenie różniczkowalna w żadnym punkcie. Podobnie $z\mapsto|z|^2$ ma rzeczywistą pochodną $h\mapsto2\operatorname{Re}(h\overline z)$, nie ma zaś zespolonej pochodnej Frécheta w punktach $z\ne0$. W punkcie $z=0$ zespolona pochodna Frécheta funkcji $z\mapsto|z|^2$ istnieje i wynosi zero, gdyż $|h|^2=o(|h|)$.
\end{przyklad}

\begin{uwaga}[Poziomy ogólności]
Związki między kierunkową a pełną różniczkowalnością, normą, topologią operatorową i różniczkowalnością zespoloną wykorzystują normowalność dziedziny. Istnieją pojęcia różniczkowania na ogólnych przestrzeniach lokalnie wypukłych (m.in. różniczkowalność Michala--Bastiani i rachunek wygodny), ale wymagają odrębnego aparatu. Nie utożsamiamy ich bez dodatkowych definicji z rachunkiem Frécheta w przestrzeniach Banacha.
\end{uwaga}
